\documentclass{oci}
\usepackage[utf8]{inputenc}
\usepackage{lipsum}
\usepackage{tikz}
\usetikzlibrary{calc}
\usepackage{tabularx}

\title{Mega Woman}

\newcommand{\megawoman}[2]{%
\node (megawoman) at ($(#1, 0)+(0.5, 0.21)$) {\scalebox{0.2}{\includegraphics{megawoman}}};
\path[->, darkgray] ($(megawoman)+(.07, 0)$) edge node[above,yshift=-2] {\scalebox{0.65}{$#2$}} ($(megawoman) + (0.6, 0)$)
}

\newcommand{\robot}[4]{%
\node (#1) at ($(#2, 0)+(0.5, 0.7)$) {\scalebox{0.1}{\includegraphics{robot}}};
\path[->, darkgray] ($(#1)+#3*(.08, 0)$) edge node[above,yshift=-2] {\scalebox{0.65}{$#4$}} ($(#1) + #3*(0.6, 0)$)
}
\newcommand{\grid}[1]{%
\draw[step=1.0,gray,thin] (0,0) grid ($(#1, 0)+(1,1)$);
\foreach \i in {0,...,#1}{
  \node at (\i+0.5,-0.2) {\scalebox{0.6}{$\i$}};
}
}
\newcommand{\bomb}[1]{%
\node at ($(#1, 0)+(0.5, 0.38)$) {\scalebox{0.2}{\includegraphics{bomb}}}
}
\newcommand*{\tabbox}[2][t]{\vspace{-10pt}\parbox[#1][4.3\baselineskip]{20em}{\strut#2\strut}}
\newcolumntype{A}{>{\footnotesize}m{.5\textwidth} }
\newcolumntype{B}{>{\vspace{10pt}}m{.45\textwidth} }

\begin{document}
\begin{problemDescription}
La reconocida empresa de videojuegos Jockeycom está desarrollando una nueva
entrega de su aclamada franquicia Megawoman.
En esta nueva versión, la heroína, Mega Woman, debe atravesar un escenario
plagado de robots bombarderos y así poder salvar a su amigo Uno.

El escenario está compuesto de $N$ casillas numeradas de 0 a $N-1$ y contiene
$R$ robots numerados de 0 a $R-1$.
Mega Woman corre siempre hacia la derecha a una velocidad constante de $v$
casillas por segundo.
Cada robot vuela por el escenario en una sola dirección, ya sea izquierda o
derecha, a una velocidad constante $u_i$.
El objetivo de Ultra Woman es escapar del escenario lográndolo cuando su
posición está más a la derecha que la última casilla.

Al comienzo de cada segundo los robots lanzan una bomba en la casilla en la que
están.
Si en ese momento Ultra Woman se encuentra dentro de la misma casilla será
alcanzada por la bomba y perderá el juego.
Los robots solo lanzan una bomba al inicio de cada segundo y, por lo tanto,
Ultra Woman puede pasar por debajo de ellos si lo hace en el momento indicado.

Dada la posición inicial de Ultra Woman y la de cada uno de los robots, a los
desarrolladores de Jockeycom les gustaría determinar si es posible para Ultra
Woman escapar del escenario sin ser alcanzada por una bomba.

La siguiente tabla muestra segundo a segundo un ejemplo de un escenario donde
Ultra Woman puede escapar del escenario sin ser alcanzada por una bomba.
La flechas indican la velocidad de cada uno de los personajes.

\begin{tabular}{|A|B|}
\hline
En el momento inicial los robots lanzan sus bombas, pero Ultra Woman no
es alcanzada por ninguna.
&
\begin{tikzpicture}[scale=1.2]
\megawoman{2}{2};
\robot{r1}{0}{-1}{1};
\bomb{0};
\robot{r2}{1}{1}{4};
\bomb{1};
\robot{r3}{5}{-1}{2};
\bomb{5};
\grid{5};
\end{tikzpicture}
\\
\hline
Transcurrido un segundo, tanto Ultra Woman como los robots se han movido la
cantidad de casillas correspondientes a su velocidad.
Adicionalmente todos los robots lanzan su bomba, pero Ultra Woman no es
alcanzada por ninguna.
&
\begin{tikzpicture}[scale=1.2]
\megawoman{4}{2};
\robot{r2}{5}{1}{4};
\bomb{5};
\robot{r3}{3}{-1}{2};
\bomb{3};
\grid{5};
\end{tikzpicture}
\\
\hline
En el segundo tres Ultra Woman logra escapar del escenario y por lo tanto gana el juego.
En este momento no importa las bombas que lancen los robots pues Ultra Woman ya se
encuentra fuera del escenario.
&
\begin{tikzpicture}[scale=1.2]
\robot{r3}{1}{-1}{2};
\grid{5};
\end{tikzpicture}
\\
\hline
\end{tabular}

\end{problemDescription}

\clearpage
\begin{inputDescription}
  La entrada está descrita en varias líneas.
  La primera contiene un entero $N$ ($N>0$), correspondiente a la cantidad de
  casillas en el escenario.
  La siguiente línea contiene dos enteros $w$ y $v$ ($0\leq w< N, 0<v\leq 10^9$)
  correspondientes a la casilla donde se encuentra inicialmente Ultra Woman y a la
  velocidad con la que se mueve.
  Posteriormente viene una línea con un entero $R$ ($R>0$)indicando la cantidad de
  robots en el escenario.
  A continuación siguen $R$ filas describiendo cada uno de los robots.
  Cada fila contiene dos enteros $r_i$ y $u_i$ ($0\leq r_i < N, |u_i|< 10^9$),
  correspondientes respectivamente a la posición inicial y a la velocidad del
  robot.
  La velocidad será positiva si el robot se mueve hacia la derecha y negativa si
  se mueve hacia la izquierda.
\end{inputDescription}

\begin{outputDescription}
La salida debe corresponder a una única línea conteniendo la palabra \texttt{SI}
en el caso que sea posible para Ultra Woman escapar del escenario o \texttt{NO}
en caso contrario.
\end{outputDescription}

\begin{scoreDescription}
  \score{20} Se probarán varios casos donde la velocidad de Ultra Woman es 1,
  todos los robots se mueven hacia la derecha con velocidad 2 y tanto la
  cantidad de casillas como de robots es menor o igual que $10^5$ ($u_i=2, v=1,
  R\leq 10^5, N\leq 10^5$).
  \score{20} Se probarán varios casos donde la velocidad de Ultra Woman es 1,
  todos los robots se mueven hacia la izquierda con velocidad 1 y tanto la cantidad
  de casillas como de robots es menor que $10^5$ ($u_i=1, v=1, R\leq 10^5,
  N\leq 10^5$).
  \score{25} Se probarán varios casos en que el número de robots es menor que
  50, el número de casillas es menor o igual que $10^5$ y sin restricciones sobre las
  velocidades ($R\leq 50, N\leq 10^5$).
  \score{35} Se probarán varios donde el número de robots es menor o igual que
  $10^5$, la cantidad de casillas es menor o igual que $10^{12}$ y sin
  restricciones sobre las velocidades ($R\leq 10^5, N \leq10^{12}$).
\end{scoreDescription}

\begin{sampleDescription}
\sampleIO{sample-1}
\sampleIO{sample-2}
\sampleIO{sample-3}
\end{sampleDescription}

\end{document}